\subsection{全章总链：从统一上位逻辑到物理时空骨架}\label{subsec:physical-spacetime-skeleton-grand-chain}

本章至此已经完成了一件在全文结构上极其关键的工作：它没有另起一套平行的物理理论，也没有用一组新的形而上学语言去覆盖前文，而是把第 \ref{sec:logic-expansion-chain} 章统一上位逻辑的 forcing 母系统，沿着观察者、区域、时间、因果、资源与谱这六条解释轴，逐步读成了一套可被物理化理解的最小时空骨架。因而，本章真正建立的，不是“从数学到物理”的跳跃，而是“从同一 forcing 骨架的不同实例层到其物理骨架化解释”的总链。

若要把这条总链压缩到最简形式，它应写成
\[
\text{统一上位逻辑}
\Longrightarrow
\text{对象层实例}
\Longrightarrow
\text{值层实例}
\Longrightarrow
\text{统一投影语法}
\Longrightarrow
\text{观察者—事件—区域骨架}
\Longrightarrow
\text{时间/因果投影}
\Longrightarrow
\text{资源几何}
\Longrightarrow
\text{谱与闭环边界不变量}.
\]

这条链中的每一箭头，都不是一种比喻，而是前文已经完成的结构提升。第 \ref{sec:logic-expansion-chain} 章给出的是全文唯一的统一上位逻辑：状态 forcing、扩张链、局部可见性、细化与保守扩张。第 \ref{sec:recursive-addressing} 章把它实例化为对象层：地址、局部截面、\(\mathsf{NULL}\)、局部可读性与决定包络。第 \ref{sec:emergent-arithmetic} 章在对象已经进入决定值域之后，把它继续实例化为值层：数值化、值保持重写、规范横截面、寄存器与账本。第 \ref{sec:pom} 章再把对象与值压缩成统一投影语法，使门、词、primitive、商余读出与证书演算进入同一高阶语言。到第 \ref{sec:physical-spacetime-skeleton} 章，这些已经存在的层不再被重新发明，而被重新阅读：观察者被读成状态纤维，事件被读成新决定，区域被读成共同支撑与局部共识，时间被读成决定包络增长的投影，因果被读成可容许精化链的偏序闭包，资源被读成这些精化链上的方向化长度，谱则被读成闭环精化的边界不变量。

因此，本章从头到尾所坚持的并不是“物理先于逻辑”，而是恰恰相反：
\[
\text{物理型时空是统一 forcing 骨架在足够丰富的实例层上所涌现出来的一种解释结构。}
\]
换言之，本章并没有把外部的时空观硬性嵌入到论文内部，而是证明：只要局部对象化、局部可读性、延迟分类、值保持、投影闭环与资源代价这些结构已经成立，那么观察者、时间、空间、因果与谱就都可以从同一骨架中导出。时空不是额外的背景，而是 forcing 状态空间的高阶组织形式。

为了把这一点说得更清楚，可以把本章前面几节的核心命题重新编号成一条“逻辑—物理”对应链。

首先，是观察者纤维化命题。它说明观察者不是一个站在系统外部的主体，而只是状态空间上的一个纤维索引。换言之，“谁在看”被还原成“哪一类状态共享同一不可区分边界”。从这个角度看，观察者并不是往系统里再加一个本体，而是对状态进行一次新的分层。

其次，是事件局部化命题。它说明事件不是坐标系中的点，而是某个先前未决定命题在某个局部区域中首次进入决定包络的那一跳。因此，事件的原初形式不是点粒子，而是“新决定的局部发生”。这一步把事件从几何点重新改写成了 forcing 中的结构突变。

再次，是空间共识命题。它说明空间不是先验几何容器，而是由共同支撑、局部比较与覆盖下的区域共识所涌现出来的最小区域系统。整体区域不是先天地在那里，而是只有当局部状态足够兼容、能够唯一拼接时，整体区域才真正作为一个可比较对象出现。

接着，是时间投影命题。它说明时间不是一个独立于状态精化之外的外部钟，而是决定包络增长的投影；同位层上的重写不构成新的时间，而真正的时间推进总对应某种新决定的出现。因此，时间箭头并不来自神秘流动，而来自决定包络的严格扩张。

随后，是因果偏序命题。它说明因果不是外加到时间之上的第二结构，而是允许更新链的偏序闭包。因果回答的不是“谁更晚”，而是“谁由谁而来”；它比时间更细，因为并不是所有较晚状态都属于某个状态的因果未来，只有那些能经由局域、可容许、可审计更新链生成的状态才属于真正的因果未来。

然后，是资源几何命题。它说明资源既不是代价表，也不是惩罚项，而是状态精化链、区域拼接和观察者同步在实现层上的方向化长度。没有资源，时间、因果与空间都只能是纯形式关系；有了资源，它们才第一次有了“实现难度”“稳定程度”“同步开销”和“共识距离”的区别。

最后，是谱边界命题。它说明谱不是附会在系统外部的装饰，而是闭环精化、时间迹与空间传播核在边界层上的压缩不变量。primitive、\(\zeta\)、压力、回返迹、传播核谱与障碍模式，全部应被读成同一个 forcing 骨架在循环结构上的不同投影。也正是在这个意义上，谱不是附属对象，而是时空骨架在闭环层的边界读法。

于是，第 \ref{sec:physical-spacetime-skeleton} 章最终给出的，并不是一份“物理解释清单”，而是一套从上位逻辑到物理骨架的生成次序。这个次序可以再次压缩成下面这组更精确的链条：
\[
(\Gamma_p,R_p)
\Longrightarrow
\widetilde p=(\Gamma_p,R_p,i,t,U,\ldots)
\Longrightarrow
[p]_i
\Longrightarrow
\mathrm{Dec}_{\mathcal O}(p)
\Longrightarrow
t_{\mathcal O}(p)
\Longrightarrow
\preceq_c
\Longrightarrow
d_{\mathcal R}
\Longrightarrow
\mathrm{Spec}.
\]

这里的每一个符号都代表一个真正的层次跳跃：
\begin{itemize}
  \item \((\Gamma_p,R_p)\) 是第 \ref{sec:logic-expansion-chain} 章中的 forcing 原状态；
  \item \(\widetilde p\) 是时空纤维化之后的局部状态；
  \item \([p]_i\) 是观察者纤维中的局部视界；
  \item \(\mathrm{Dec}_{\mathcal O}(p)\) 是可观测片段上的决定包络；
  \item \(t_{\mathcal O}(p)\) 是由决定包络增长所定义的时间投影；
  \item \(\preceq_c\) 是由可容许精化链导出的因果偏序；
  \item \(d_{\mathcal R}\) 是这些精化链上的资源长度；
  \item \(\mathrm{Spec}\) 是这些长度与闭环统计压缩成的谱边界不变量。
\end{itemize}

因此，若要问“物理型时空究竟是什么”，本章给出的最终回答不是“一个带度量的流形”或“一个事件集合加因果锥”，而是
\[
\text{物理型时空是统一上位逻辑 forcing 状态在观察者纤维化、区域共识、决定包络投影、可容许精化、资源长度与闭环谱压缩之下所涌现出的综合结构。}
\]

这个定义的力量在于，它不要求一开始就接受某一种几何本体，也不要求一开始就站在某一种物理理论的内部。它只要求承认：若一个系统已经具备局部对象化、局部可读性、状态精化、局部拼接、值保持、投影闭环与资源代价，那么“观察者—空间—时间—因果—谱”就不再是五个互不相干的额外对象，而是同一 forcing 母系统的五个边界投影。

由此也可以反过来说明本章的边界。它并没有证明“某一具体物理理论必然由这些结构给出”，也没有证明“任何具备这些结构的对象都自动成为现实物理世界”。它证明的只是：
\[
\text{在论文当前已经建立起来的 forcing 骨架内部，存在一条足够强、足够完整、足以被物理化解释的时空骨架链。}
\]
换言之，本章完成的是“可物理化性”的证明，而不是“物理真实性”的宣告。

最后，本章的总链还应当被回指到全文整体结构中去。第 \ref{sec:logic-expansion-chain} 章是母系统；第 \ref{sec:recursive-addressing} 章是对象层实例；第 \ref{sec:emergent-arithmetic} 章是值层与实现层实例；第 \ref{sec:pom} 章是统一投影语法层；第 \ref{sec:physical-spacetime-skeleton} 章则把这四层强行压到一条“观察者—时空—资源—谱”的解释链上。因此，从全书结构上说，第 \ref{sec:physical-spacetime-skeleton} 章既不是附录性的哲学评论，也不是对前文技术工作的重复说明，而是一个真正的“解释封顶章”：它把前文所有强技术块收束为一个单一母系统的远端几何读法。

因此，本节可以以如下结论收尾：
\[
\text{从第 4 章到第 15 章，论文并未更换逻辑基础；变化的只是同一 forcing 骨架的解释层级。第 4 章把它写成上位逻辑，第 7–9 章把它写成对象、值与统一语法的实例，而第 15 章把它最终读成物理型时空骨架。}
\]

这条总链一旦立住，全文的逻辑层次就会变得非常清楚：不再是若干相互平行的技术章节，而是一套从 forcing 母系统逐层外展、最后在物理骨架上闭合的统一工程。
